\section{Conclusion}
\label{sec:conclusion}

Score-weighted metaheuristics are justified by a causal claim---that a
per-neighbour score, computed on the population graph or on the run's own
history, carries information about where to search---which an ablation cannot
establish, because deleting the component changes the mixing operator along
with the information it carries.

We defined two controls that keep the mechanism running at full strength and
destroy only the information it is claimed to use. Permuting a score across
nodes yields an operator equal to uniform mixing in expectation that retains
the original's weight heterogeneity in every realisation
(Proposition~\ref{prop:mean})---exactly the separation an ablation fails to
make---while a leverage bound (Proposition~\ref{prop:leverage}) states when
such a comparison is informative at all.

Applied to a published betweenness-weighted optimizer, the controls give a
clear answer. The mechanism is active, displacing each update by a quarter to a
third of the population's spread. Its assignment of scores to graph positions
is not: across $104$ function comparisons on two suites, three dimensions and
five budgets, permutation changes nothing that survives correction, and at
suite level the effect is equivalent to zero within a margin four times
stricter than negligibility. Unequal weights help slightly, on $3$ of $28$
functions; which neighbour receives the large weight matters on none.

A null result of this kind is only as good as the test's ability to return a
non-null one, so we calibrated the instrument on four further mechanisms,
unchanged in suite, sample size and statistic. Within the same optimizer it
returns an intermediate, class-dependent effect on the fitness-influence
weighting and a clear one on adaptive density selection ($\bar{\delta}=-0.101$,
$p=1.5\times10^{-5}$). On a second published optimizer, the fully informed
particle swarm, it returns a uniformly harmful effect for a synthetic
degree weighting of our own construction, and for the published quality
weighting something else again: per-function effects reaching $|\delta|=1$ in
both directions, significant on $15$ of $28$ functions, that cancel to
$\bar{\delta}=-0.090$ in the mean. Five distinguishable outcomes---inert by
construction, uninformative, sign-varying, harmful, beneficial---on one suite
with one statistic. A null returned by a test with that range is a statement
about the mechanism.

That fourth case also settles a question of reporting: centrality and quality
weighting fail the same suite-level test for opposite reasons---nothing happens
anywhere, against a great deal happening in both directions---while our own
adaptive scheme shows the mirror failure, an effect unambiguous in aggregate
that per-function correction hides. Neither view is sufficient and we report
both.

What carries information in this algorithm, then, is not the scores computed on
the graph but the graph's coarse density: consequential in $30$ of $49$
configurations, with an optimum that is problem- and dimension-dependent and
resists prediction from a ruggedness descriptor. We built an adaptive scheme on
that basis and subjected it to the second control, which it passes.

The instrument costs one arm and no additional budget, and we would encourage
its use wherever a structural or adaptive claim is made, under two conditions
our own analysis makes concrete: report the displacement, so that a control on
an inert mechanism is not mistaken for evidence, and test equivalence
explicitly with adequate power rather than resting on a non-significant
difference. The controls apply unchanged to any score-proportional weighting,
so the audit can be widened directly; the constructive extension is to ask
whether a graph derived from the landscape, rather than fixed in advance, would
make a structural statistic informative. The prediction these controls can test
is that on such a graph permutation would hurt.
